\section{1868 Beltrami, Cayley i Klein} % lang-pl
\section{1868 Beltrami, Cayley e Klein} % lang-it
\label{sekcja_beltrami_cayley_klein}
Opowiemy tutaj o modelu Beltramiego-Kleina dla geometrii hiperbolicznej. % lang-pl
Parleremo qui del modello di Beltrami-Klein per la geometria iperbolica. % lang-it
Zaczniemy od definicji, potem przedstawimy jego historię. % lang-pl
Cominceremo dalla definizione, poi ne presenteremo la storia. % lang-it

In two dimensions the Beltrami–Klein model is called the Klein disk model. It is a disk and the inside of the disk is a model of the entire hyperbolic plane. Lines in this model are represented by chords of the boundary circle (also called the absolute). The points on the boundary circle are called ideal points; although well defined, they do not belong to the hyperbolic plane. Points outside the disk do not belong to the hyperbolic plane either, and they are sometimes called ultra ideal points.



W dwóch wymiarach model Beltramiego-Kleina nazywa się raczej modelem dysku Kleina. % lang-pl
In due dimensioni il modello di Beltrami-Klein si chiama piuttosto modello del disco di Klein. % lang-it
Wnętrze dysku realizuje w nim płaszczyznę hiperboliczną. % lang-pl
L'interno del disco vi realizza il piano iperbolico. % lang-it
Prostymi w tym modelu są cięciwy brzegu dysku; punkty na brzegu nazywa się doskonałymi i nie należą do płaszczyzny. % lang-pl
Le rette in questo modello sono le corde del bordo del disco; i punti sul bordo si dicono ideali e non appartengono al piano. % lang-it
Model ten zniekształca kąty oraz okręgi, których środki nie leżą na środku wyjściowego dysku. % lang-pl
Questo modello deforma gli angoli e le circonferenze i cui centri non coincidono con il centro del disco di partenza. % lang-it
Dwie proste (cięciwy) są prostopadłe, jeśli po przedłużeniu poza dysk, każda przechodzi przez biegun drugiej. % lang-pl
Due rette (corde) sono perpendicolari se, prolungate oltre il disco, ciascuna passa per il polo dell'altra. % lang-it

% Both the Poincaré disk model and the Klein disk model are models of the hyperbolic plane. An advantage of the Poincaré disk model is that it is conformal (circles and angles are not distorted); a disadvantage is that straight lines project to circular arcs orthogonal to the boundary circle of the disk.

W 1859 roku Arthur Cayley użyje dwustosunkowej definicji kąta (od Laguerre'a), by pokazać, że wewnątrz geometrii euklidesowej można zdefiniować geometrię rzutową. % lang-pl
Nel 1859 Arthur Cayley utilizzerà la definizione di angolo tramite rapporto anarmonico (dovuta a Laguerre) per mostrare che all'interno della geometria euclidea si può definire la geometria proiettiva. % lang-it
W 1869 roku młody Felix Klein zapozna się z dorobkiem Cayleya % lang-pl
Nel 1869 il giovane Felix Klein conoscerà l'opera di Cayley. % lang-it
W 1868 roku Eugenio Beltrami, matematyk włoski, opublikuje dwie prace o tym, że geometria hiperboliczna oraz euklidesowa są równoważnie niesprzeczne: jedna jest niesprzeczna dokładnie wtedy, kiedy druga ma tę cechę. % lang-pl
Nel 1868 Eugenio Beltrami, matematico italiano, pubblicherà due lavori per mostrare che la geometria iperbolica e quella euclidea sono equiconsistenti: l'una è coerente se e solo se lo è l'altra. % lang-it
Najpierw pojawi się \emph{Saggio di interpretazione della geometria non-euclidea} (przypadek pseudosfery, jej geodezyjne badali wcześniej Ferdinand Minding i Delfino Codazzi); później \emph{Teoria fondamentale degli spazii di curvatura costante} (ogólne rozumowanie dla dodwolnej liczby wymiarów). % lang-pl
Dapprima apparirà il \emph{Saggio di interpretazione della geometria non-euclidea} (il caso della pseudosfera, le cui geodetiche avevano studiato prima Ferdinand Minding e Delfino Codazzi); poi la \emph{Teoria fondamentale degli spazii di curvatura costante} (ragionamento generale per un numero qualsiasi di dimensioni). % lang-it
Prace Beltramiego nie zdobędą rozgłosu. % lang-pl
I lavori di Beltrami non otterranno grande risonanza. % lang-it



The papers of Beltrami remained little noticed until recently and the model was named after Klein ("The Klein disk model"). In 1859 Arthur Cayley used the cross-ratio definition of angle due to Laguerre to show how Euclidean geometry could be defined using projective geometry.[4] His definition of distance later became known as the Cayley metric.



powoła się przy tym na wykład Riemanna z 1854 roku. % lang-pl
Si richiamerà a tal fine alla conferenza di Riemann del 1854. % lang-it
\index[persons]{Hoüel, Jules}%
\index[persons]{Minding, Ferdinand}%
\index[persons]{Codazzi, Delfino}%

Kolejnym osiągnięciem będzie model Beltramiego-Kleina dla geometrii hiperbolicznej, w którym pracujemy z punktami wewnątrz otwartego koła. % lang-pl
Proste w tym modelu to cięciwy łączące punkty z brzegu koła. % lang-pl
Wspomni o nim Eugenio Beltrami w 1868 roku, najpierw dla $n = 2$, a potem dla dowolnej liczby wymiarów. % lang-pl
Jego prace dowiodą, że geometria hiperboliczna jest równoniesprzeczna z euklidesową. % lang-pl
Un ulteriore progresso sarà il modello di Beltrami-Klein della geometria iperbolica, nel quale si lavora con i punti all'interno di un disco aperto. % lang-it
Le rette di questo modello sono le corde che uniscono punti del bordo del disco. % lang-it
Eugenio Beltrami lo presenterà nel 1868, dapprima per $n = 2$ e successivamente per un numero arbitrario di dimensioni. % lang-it
Le sue opere proveranno che la geometria iperbolica è equiconsistente con quella euclidea. % lang-it

% memoirs
Wcześniej, bo w 1859 roku, Arthur Cayley użyje dwustosunkowej definicji kąta (od Laguerre'a) by pokazać, że geometria euklidesowa może być zdefiniowana przy pomocy geometrii rzutowej. % lang-pl
In precedenza, nel 1859, Arthur Cayley utilizzerà la definizione di angolo tramite rapporto anarmonico (dovuta a Laguerre) per mostrare che la geometria euclidea può essere definita mediante la geometria proiettiva. % lang-it
\index[persons]{Cayley, Arthur}%
\index[persons]{Laguerre, Edmond}%
\index[persons]{Sylvester, James}%
Ten wynik ukaże się w \emph{A sixth memoir upon quantics}; odległość, którą tam zdefiniuje za pomocą logarytmu dwustosunku, zostanie później nazwana metryką Cayleya. % lang-pl
Cayley nazwie ponadto zbiór punktów idealnych płaszczyzny Łobaczewskiego jej \emph{absolutem} \cite[s. 307]{eves1_1972} (zobacz sekcję \ref{sekcja_punkty_idealne}). % lang-pl
Questo risultato apparirà in \emph{A sixth memoir upon quantics}; la distanza che vi definirà tramite il logaritmo del birapporto sarà poi chiamata metrica di Cayley. % lang-it
Cayley chiamerà inoltre il luogo dei punti ideali del piano di Lobachevskij il suo \emph{assoluto} \cite[p. 307]{eves1_1972} (si veda la sezione \ref{sekcja_punkty_idealne}). % lang-it

Klein zapozna się z pracami Beltramiego i Cayleya; w 1870 roku, jako dwudziestolatek, weźmie udział w seminarium Weierstrassa, gdzie omówi pomysły Cayleya i zapyta: czy między pomysłami Cayleya oraz Łobaczewskiego istnieje powiązanie? % lang-pl
Usłyszy w odpowiedzi, że oba systemy dzieli przepaść pojęciowa. % lang-pl
Klein studierà i lavori di Beltrami e Cayley; nel 1870, appena ventenne, parteciperà al seminario di Weierstrass, dove discuterà le idee di Cayley e chiederà: esiste una relazione tra le idee di Cayley e quelle di Lobachevskij? % lang-it
Riceverà in risposta che i due sistemi sono concettualmente molto lontani. % lang-it
\index[persons]{Klein, Felix}%
\index[persons]{Weierstrass, Karl}%
Klein nie zrezygnuje: w 1871 roku opublikuje pracę \emph{Über die sogenannte Nicht-Euklidische Geometrie} i pokaże w niej, że metody Cayleya dają rzutowy model geometrii nieeuklidesowej. % lang-pl
A rok później, w 1872, wygłosi program erlangeński. % lang-pl
Choć Beltrami był pierwszy, model przejdzie do historii jako model Kleina (a sprawiedliwiej -- Cayleya--Kleina); prace Beltramiego z 1868 roku pozostaną w dużej mierze niezauważone. % lang-pl
Klein non si arrenderà: nel 1871 pubblicherà il lavoro \emph{Über die sogenannte Nicht-Euklidische Geometrie} e vi mostrerà che i metodi di Cayley forniscono un modello proiettivo della geometria non euclidea. % lang-it
E un anno dopo, nel 1872, terrà il programma di Erlangen. % lang-it
Anche se Beltrami era stato il primo, il modello passerà alla storia come modello di Klein (più giustamente: di Cayley--Klein); i lavori di Beltrami del 1868 rimarranno in gran parte inosservati. % lang-it
\index{model!Beltramiego-Kleina} % lang-pl
\index{modello!di Beltrami-Klein} % lang-it

\begin{definition}
[model Beltramiego-Kleina] % lang-pl
[modello di Beltrami-Klein] % lang-it
	Ustalmy okrąg $\Sigma$. % lang-pl
	Punktami są punkty wnętrza $\Sigma$, prostymi -- cięciwy $\Sigma$, a długość odcinka $AB$ to $\log(AB, TS)$, gdzie $S$, $T$ są końcami cięciwy zawierającej $AB$, uporządkowanymi tak, że $A$ leży między $S$ i $B$. % lang-pl
	Odcinki (kąty) są przystające, gdy jeden przechodzi na drugi przez rzutowość przeprowadzającą $\Sigma$ i jej wnętrze na siebie. % lang-pl
	Fissiamo una circonferenza $\Sigma$. % lang-it
	I punti sono i punti dell'interno di $\Sigma$, le rette sono le corde di $\Sigma$, e la lunghezza del segmento $AB$ è $\log(AB, TS)$, dove $S$, $T$ sono gli estremi della corda contenente $AB$, ordinati in modo che $A$ stia tra $S$ e $B$. % lang-it
	Segmenti (angoli) sono congruenti quando uno si trasforma nell'altro tramite una proiettività che manda $\Sigma$ e il suo interno in sé. % lang-it
\end{definition}

% Eves \cite[s. 352, zad. 6--7]{eves1_1972}, gdzie przypisuje ten model Kleinowi. % lang-pl
% Eves \cite[p. 352, es. 6--7]{eves1_1972}, dove attribuisce questo modello a Klein. % lang-it
Jarosław Górnicki przedstawi w $\Delta_{20}^{8}$ inny model płaszczyzny hiperbolicznej: wnętrze koła, w którym proste to średnice lub łuki okręgów przecinających brzeg pod kątem prostym, a kąty są wiernie przenoszone z płaszczyzny euklidesowej. % lang-pl
Przypisze go Beltramiemu (1868), dodając, że 14 lat później podał go także Henri Poincaré; w tym modelu przez punkt poza prostą przechodzą dwie (w istocie nieskończenie wiele) proste jej nieprzecinające, co posłuży do uzasadnienia niezależności piątego aksjomatu od czterech pozostałych. % lang-pl
Jarosław Górnicki presenterà in $\Delta_{20}^{8}$ un altro modello del piano iperbolico: l'interno di un disco, in cui le rette sono diametri o archi di circonferenze che incontrano il bordo ad angolo retto, e gli angoli sono riportati fedelmente dal piano euclideo. % lang-it
Lo attribuirà a Beltrami (1868), aggiungendo che 14 anni dopo lo diede anche Henri Poincaré; in questo modello per un punto fuori da una retta passano due (in realtà infinite) rette che non la intersecano, il che servirà a giustificare l'indipendenza del quinto assioma dagli altri quattro. % lang-it

Beltrami odegra jeszcze jedną rolę w naszej opowieści: w 1889 roku przypomni światu zapomnianą książkę Saccheriego (zobacz sekcję \ref{sekcja_saccheri}). % lang-pl
Beltrami avrà ancora un altro ruolo nella nostra storia: nel 1889 richiamerà l'attenzione del mondo sul dimenticato libro di Saccheri (si veda la sezione \ref{sekcja_saccheri}). % lang-it
% Źródła: Hartshorne s. 334, 355--366 (§37, §39); Eves s. 347--353 (§8.5), 353 nn. (§8.6); Greenberg 2010, s. 209, 213--214; Wikipedia: Poincaré disk model, Henri Poincaré, Felix Klein.
% https://en.wikipedia.org/wiki/Poincar%C3%A9_disk_model
% Uwaga: Wikipedia przypisuje metrykę dysku już Beltramiemu (1868); nazwa pochodzi od Poincarego (ok. 1882).

